\begin{abstract}
多核 NPU 执行一张计算图，需要决定每个操作属于哪个子图、每个子图放在哪个核心、同核子图的先后顺序。三项决策相互牵制：切得细能并行，但跨子图数据要经共享 DDR 往返并付出同步延迟；切得粗能复用核内缓存，但并行度下降，大子图还会引发换入换出。本文以官方评估程序为唯一评价标准，先建立带无环约束的切图—分核—排序联合模型并证明其强 NP 难，再用解析下界、零通信表调度和含同步延迟的可达参考衡量结果离可达水平还有多远。测试集的依赖结构几乎不限制并行（4 核理想并行效率均值 \V{ideal_k4_eff_mean}），差距主要来自通信与同步。按最大连通分量的计算量占比，100 个用例可分为单一巨块、主块加碎片和多分量三类，适合它们的切法各不相同，这决定了本文“多种构造、逐用例选优”的求解框架。

\textbf{问题一（场景 A）。}每个子图是一个 Task，跨子图数据必须经 DDR 中转并伴随 100/1000 周期的激活等待。我们推导了精确的边界搬运量公式，并证明固定延迟沿子图依赖路径线性累积，从而得到“少跨核、少在依赖链上切分”的设计原则（用扫描候选检验，代价取决于路径深度而非子图总数）。据此构造三种算法，每种都给出无环性证明：分量装箱与拓扑区间切分（适合多分量图）、层带$\times$核划分（适合单一巨块图）、通信感知表调度加相位构造（子图数只由路径上的换核次数决定）。三种算法在参数网格上逐用例选优后，再依次做成组合并、局部搜索和核数补齐。2--5 核平均加速比依次为 \V{p1_k2_final_mean}、\V{p1_k3_final_mean}、\V{p1_k4_final_mean}、\V{p1_k5_final_mean}，5 核比朴素分量装箱（\V{p1_k5_naive_mean}）提高 \V{p1_k5_naive_gain_pct}，对含同步可达参考的平均达成率为 \V{p1_k5_achieve_mean}。各步骤中，逐用例选优和成组合并贡献了主要提升，局部搜索平均缩短 \V{ls_p1_mean_red_pct}，只在个别用例上不可替代。

\textbf{问题二（场景 B）。}同核子图合并为一个 Task，数据可在 L1/UB 中复用。我们证明同一方案下场景 B 的边界搬运量不超过场景 A，时间代价由按子图累积的激活等待变为按换核次数累积的 500 周期同步延迟；同时加入缓存驻留约束，刻画跨子图驻留数据引发的换出。模型预测场景 B 适合细分，扫描数据验证了这一点：5 核最终方案的子图数中位数为 \V{p2_k5_nsg_median}，扫描候选的质量随子图数增加单调变好。2--5 核平均加速比依次为 \V{p2_k2_final_mean}、\V{p2_k3_final_mean}、\V{p2_k4_final_mean}、\V{p2_k5_final_mean}，对可达参考的达成率为 \V{p2_k5_achieve_mean}（5 核）；单一巨块类 5 核平均加速比为 \V{p2_k5_famI_mean}。

\textbf{问题三（共享只读 L2 Cache）。}我们写出命中的时序条件与容量条件，把未命中分解为首次、并发、容量三类，并用 DDR 忙碌占比刻画命中能否转化为时间收益。问题三与问题二同属场景 B，候选方案在两问之间交叉复评，并在缓存评估下做成组合并和局部搜索。1--5 核只读 Cache 配置的平均加速比为 \V{p3_k1_cache_vs_single}、\V{p3_k2_cache_vs_single}、\V{p3_k3_cache_vs_single}、\V{p3_k4_cache_vs_single}、\V{p3_k5_cache_vs_single}，无 L2 为 \V{p3_k1_noL2_vs_single}--\V{p3_k5_noL2_vs_single}；同核数下只读 Cache 相对无 L2 的加速比从 1 核的 \V{p3_k1_best_mean} 升至 5 核的 \V{p3_k5_best_mean}，单例最高 \V{p3_k5_best_max}，平均命中率从 \V{p3_k1_hit_mean} 升至 \V{p3_k5_hit_mean}。4 核缺失分解中首次缺失占 \V{fig_miss_compulsory}、并发缺失占 \V{fig_miss_concurrent}。容量与带宽的敏感性实验表明，官方配置附近的瓶颈不在缓存规格，而在读取时序；对照实验中把缓存带宽从 250 降到 60 bytes/cycle，锁步读取被打散，命中率反而从 \V{fig_lockstep_hit250_pct} 升至 \V{fig_lockstep_hit60_pct}。

全部结果由官方评估程序终版复评得到，1{,}400 个主结果与 500 组配对逐位一致；三个问题在 2--5 核下均满足核数单调；参数扰动 $\pm20\%$ 的鲁棒性实验中，本文方案相对朴素方案的优势保持稳定。

\keywords{多核调度\quad 带无环约束的图划分\quad 相位构造\quad 只读缓存\quad 读取时序}
\end{abstract}
